% ================================================================= BASELINE TỰ CÀI LẠI
\subsection{Baseline cổ điển tự cài lại}
\label{sec:baselines}

Bản đề cương trước thừa nhận một điểm yếu: mọi so sánh với y văn đều là so với \textit{con số đã công
bố}, không phải với bản tự chạy. Nhóm đã khắc phục. Ba baseline không dùng học máy được cài lại từ đầu
trong \texttt{baselines/ts\_baseline.py} và chấm trên \textbf{cùng dữ liệu, cùng giao thức, cùng hàm khớp
nhịp} với mô hình của nhóm.

\begin{table}[H]\centering\small
\caption{Ba baseline cổ điển. Siêu tham số của mỗi baseline được chọn \textbf{chỉ trên r01} rồi áp
nguyên cho bốn bản ghi còn lại; cột ``16 giữ lại'' loại r01 khỏi phép tính.}
\label{tab:baseline_mota}
\begin{tabular}{@{}L{2.6cm} L{6.2cm} L{5.6cm}@{}}
\toprule
\textbf{Baseline} & \textbf{Khử ECG mẹ} & \textbf{Dò QRS thai} \\
\midrule
TS & Mẫu trung bình, trừ trực tiếp & Pan--Tompkins thích nghi, trơ 250\,ms \\
TS-PCA & Mẫu từ 2 thành phần chính của các nhịp mẹ (theo Behar 2014) & Pan--Tompkins thích nghi \\
Prominence & \textbf{Đúng front-end của nhóm}: mẫu trung vị + bình phương tối thiểu từng nhịp &
\texttt{find\_peaks}, độ nhô $\geq$ trung vị $+ 3\,$MAD, cách nhau $\geq$ 250\,ms \\
\bottomrule
\end{tabular}
\end{table}

Baseline thứ ba là quan trọng nhất về mặt khoa học: nó dùng \textbf{cùng tín hiệu dư} mà mạng nhận vào,
chỉ thay mạng bằng phép lấy đỉnh đơn giản. Hiệu số giữa nó và mô hình đo \textbf{đúng phần đóng góp của
mạng nơ-ron}, tách khỏi phần đóng góp của tiền xử lý.

\begin{table}[H]\centering\small
\caption{Kết quả trên ADFECGDB, dung sai $\pm$50\,ms. Hiệu số và khoảng tin cậy tính ở \textbf{mức chủ
thể} ($n = 5$ sản phụ, trung bình 4 kênh trước). Các trị số $p$ của bản 3.2 tính trên 20 cặp (bản ghi
$\times$ kênh) là \textbf{giả lập} và đã bị gỡ --- xem \S\ref{sec:thongke}.}
\label{tab:baseline_kq}
\begin{tabular}{@{}L{2.4cm} C{2.2cm} C{1.4cm} C{1.6cm} C{1.4cm} C{2.4cm} C{2.4cm}@{}}
\toprule
\textbf{Phương pháp} & \textbf{4 kênh (20)} & \textbf{16 giữ lại} & \textbf{Kênh PSD} &
\textbf{Hiệu số} & \textbf{KTC 95\% bootstrap} & \textbf{KTC 95\% $t$} \\
\midrule
TS + PT & $78{,}96 \pm 23{,}52$ & 75,55 & 87,68 & $-18{,}49$ & $[-30{,}60; -6{,}38]$ &
\xau{$[-38{,}80; +1{,}81]$} \\
TS-PCA + PT & $91{,}05 \pm 10{,}17$ & 90,94 & 96,74 & $-6{,}40$ & $[-8{,}09; -4{,}70]$ &
\tot{$[-9{,}10; -3{,}69]$} \\
Prominence & $86{,}39 \pm 11{,}14$ & 84,99 & 92,03 & $-11{,}06$ & $[-14{,}51; -7{,}80]$ &
\tot{$[-16{,}38; -5{,}73]$} \\
\midrule
\textbf{\mh{}} & $\mathbf{97{,}45 \pm 4{,}44}$ & \textbf{96,81} & \textbf{99,21} & --- & --- & --- \\
\bottomrule
\end{tabular}
\end{table}

\begin{canhbao}[title={Ba baseline này KHÔNG cùng độ chắc --- bản 3.2 trình bày chúng như nhau}]
Với \textbf{TS-PCA} và \textbf{Prominence}, KTC 95\% của hiệu số loại trừ 0 ở \textbf{cả} bootstrap lẫn
thang $t$ --- hướng kết quả vững, chỉ là không được tuyên bố ý nghĩa thống kê từ 5 phụ nữ.

Với \textbf{TS} thì khác: KTC 95\% kiểu $t$ là $[-38{,}80;\ +1{,}81]$ --- \textbf{chứa 0}. Hiệu số
18,49 điểm hoàn toàn do biến thiên giữa 5 bản ghi, không phải một lợi thế ổn định. Đây là baseline yếu
nhất về mặt bằng chứng, dù trông là baseline nhóm thắng đậm nhất.
\end{canhbao}

Ba điều rút ra, xếp theo mức quan trọng:

\begin{enumerate}[leftmargin=1.6em,itemsep=4pt]
\item \textbf{Mạng nơ-ron đáng $+11{,}06$ điểm} so với lấy đỉnh trên cùng tín hiệu dư, KTC 95\% kiểu $t$
$[+5{,}73;\ +16{,}38]$ loại trừ 0. Đây là con số trả lời câu hỏi ``mạng có thật sự cần không'' --- và câu
trả lời là có, nhưng chỉ khi front-end đã đúng.
\item \textbf{TS-PCA là baseline cổ điển mạnh.} Trên kênh chọn mù nhãn nó đạt 96,74, chỉ kém mô hình
2,47 điểm. Điều này nhất quán với Power-MF (98,0 bằng phương pháp cổ điển) và củng cố luận điểm chính:
phần lớn hiệu năng nằm ở front-end, không ở mạng.
\item \textbf{Độ lệch chuẩn giảm hơn một nửa} (4,44 so với 10,17 của baseline tốt nhất). Mạng không chỉ
cao hơn trung bình mà còn \textit{ổn định hơn} giữa các sản phụ và kênh.
\end{enumerate}

\begin{ghichu}[title={Power-MF: bản 3.2 nói không chạy lại được --- bản 3.4 đã chạy, ở hai cấu hình}]
Bản 3.2 ghi rằng Power-MF cần MATLAB R2021a và các hàm ICA đa kênh của Varanini 2014 không có trong kho
mã công bố, nên không chạy lại được. \textbf{Phát biểu đó đã bị rút.} Power-MF chạy được trên
\textbf{GNU Octave 11.3.0} sau bảy bản vá, và được chấm bằng \textbf{chính bộ chấm của nhóm} trên
\textbf{cả 22 chủ thể}. Con số công bố cho Silesia B1 là 99,46; nhóm chạy lại được \textbf{99,40} ---
lệch 0,06 điểm.

Nhóm cài thêm \textbf{Power-MF-1ch}: cùng thuật toán, cùng front-end và bộ khử mẹ của nhóm, nhưng chỉ
\textbf{một đạo trình}. Đây là đối chứng công bằng nhất trong cả đề cương, vì mọi thứ trừ thuật toán lõi
đều giống hệt nhau. Nó đạt 86,71 --- kém \hethong{} 10,85 điểm, thua 22/22 chủ thể.

\xau{Cảnh báo:} mọi con số Power-MF của bản 3.3 (94,87 / 98,38 / 97,33 / hiệu số $+2{,}74$ và $-1{,}06$)
đứng trên một bản Octave hỏng vì lỗi cổng chuyển của nhóm, và \textbf{đã bị rút toàn bộ}. Xem hộp đầu
\S\ref{sec:powermf}.
\end{ghichu}
